% Diagram of single-slit diffraction setup with an inset closeup of the slit.
% Based on https://tikz.net/optics_diffraction/, with some cosmetic adaptations.
\begin{tikzpicture}

    % Closeup view of the single slit
    \begin{scope}[shift={(-7,0)}, scale=1.2]
        \def\L{6.15}      % distance between walls
        \def\l{4.0}       % path length
        \def\H{4.3}       % total wall height
        \def\f{0.9}       % fractional height of projection point
        \def\t{0.15}      % wall thickness
        \def\a{2.5}       % slit distance
        \def\ang{26}      % angle
        \coordinate (T) at (0,\a/2);
        \coordinate (B) at (0,-\a/2);
        \coordinate (L) at (0,0);
        \coordinate (R) at (\L,0);
        \coordinate (I) at ({\a/2/tan(\ang)},0);

        % LINES
        \draw[photon] (T) --++ (\ang:.8*\l) coordinate (PT); 
        %node[midway,below=1,above left=-2] {$r_1$};
        \draw[photon] (B) --++ (\ang:\l) coordinate (PB); %node[midway,left=2,below right=-1] {$r_2$};
        \draw[photon] (L) --++ (\ang:.9*\l) coordinate (PR);
        \draw[dashed] (T) -- (B);
        \draw[RocBlue,dashed] (T) -- ($(L)!(T)!(PR)$) coordinate (M);
        %\draw[black!60,dashed] (T) --++ (.20*\L,0) coordinate (TR);
        \draw[black!60,dashed] (L) --++ (.6*\L,0) coordinate (LR);
        %\draw[black!60,dashed] (B) --++ (.20*\L,0) coordinate (BR);

        % LINE END
        \lineend{PT}{\ang+70}
        \lineend{PR}{\ang+70}
        \lineend{PB}{\ang+70}

        % ANGLES
        \draw pic[-latex,"$\theta$",RocBlue,draw=RocBlue,angle radius=20,angle eccentricity=1.3]
          {angle = B--T--M}; %\contour{white}{}
        \draw pic[-latex,"$\theta$",RocBlue,draw=RocBlue,angle radius=24,angle eccentricity=1.2]
          {angle = LR--L--PR};
        \draw pic[-latex,"$\theta$",RocBlue,draw=RocBlue,angle radius=24,angle eccentricity=1.2]
          {angle = LR--I--PB};
        \rightAngle{T}{M}{PR}{0.3}

        % MEASURES
        \draw[latex-latex,black] (-2.1*\t,0) --++ (0,\a/2) node[midway,fill=white,inner sep=0.1] {$a/2$};
        \draw[latex-latex,black] (-2.1*\t,0) --++ (0,-\a/2) node[midway,fill=white,inner sep=0.1] {$a/2$};
        \draw[latex-latex,black] ([shift={({\ang-90}:0.1)}]L) -- ([shift={({\ang-90}:0.1)}]M)
          node[midway,left=5,below right=-1]{$\displaystyle{\frac{a}{2}\!\sin\theta}$};

        % WALL
        \fill[gray]
          (0,\a/2) rectangle (-\t,\H/2)
          (0,-\a/2) rectangle (-\t,-\H/2);
  \end{scope}

    % Single slit diagram showing the screen
    \begin{scope}
        \def\L{9}         % distance between walls
        \def\H{4.3}       % total wall height
        \def\f{0.95}      % fractional height of projection point
        \def\ang{atan((\f*\H+\a/2)/\L/2)} % theta
        \def\t{0.15}      % wall thickness
        \def\a{2.5}       % slit distance
        \coordinate (T) at (0,\a/2);
        \coordinate (B) at (0,-\a/2);
        \coordinate (L) at (0,0);
        \coordinate (R) at (\L,0);
        \coordinate (P) at (\L,\f*\H/2);
        \coordinate (M) at ($(L)!(T)!(P)$);

        % LINES
        \draw[photon] (T) -- (P); %node[midway,above=-1] {$r_1$};
        \draw[photon] (B) -- (P); %node[midway,below=3,right=6] {$r_2$}; %right=6,below right=-4
        \draw[photon] (L) -- (P);
        \draw[dashed] (T) -- (B);
        \draw[dashed,black!60] (L) -- (R);
        \draw[RocBlue,dashed] (M) -- (T);

        % ANGLES
        \draw pic[-,"\footnotesize$\theta$",RocBlue,draw=RocBlue,angle radius=25,angle eccentricity=1.2]
          {angle = B--T--M};
        \draw pic[-,"\footnotesize$\theta$",RocBlue,draw=RocBlue,angle radius=31,angle eccentricity=1.14]
          {angle = R--L--P};
        \rightAngle{T}{M}{P}{0.3}

        % MEASURES
        \draw[latex-latex,black] (0,-0.47*\H) --++ (\L,0) node[midway,fill=white,inner sep=1] {$L$};
        \draw[latex-latex,black] (-2.1*\t,0) --++ (0,\a/2) node[midway,fill=white,inner sep=0.1] {$a/2$};
        \draw[latex-latex,black] (-2.1*\t,0) --++ (0,-\a/2) node[midway,fill=white,inner sep=0.1] {$a/2$};
        \draw[|-|,black] ([shift={({\ang-100}:0.1)}]L) -- ([shift={({\ang-100}:0.1)}]M)
          node[midway,left=3,below right]{\footnotesize $\displaystyle{\frac{a}{2}\sin\theta}$};
        \draw[latex-latex,black] ([shift={(2.1*\t,0)}]P) -- ([shift={(2.1*\t,0)}]R)
          node[midway,fill=white,inner sep=1]{$y$};

        % WALL
        \fill[gray]
          (0,\a/2) rectangle (-\t,\H/2)
          (0,-\a/2) rectangle (-\t,-\H/2)
          (\L,-\H/2) rectangle (\L+\t,\H/2);
        \draw[YellowOrange!80!black, fill=white] (P) circle (0.4*\t) node[above left] {$P$};
%        \fill[YellowOrange!80!black] (P) circle (0.3*\t) node[right=1,above left=-1] {P};
    \end{scope}

    % Draw a bounding box around part of the scene:
    \draw[densely dotted] (-0.75,-2.25) rectangle (1.75,2.25);

    % Draw a bounding box around the closeup with connections to the larger
    % drawing to show the inset:
    \draw[densely dotted] (-7.75,-2.75) rectangle (-2.5,3.5);
    \draw[densely dotted] (-2.5,-2.75) -- (1.75,-2.25);
    \draw[densely dotted] (-2.5,3.5) -- (1.75,2.25);

\end{tikzpicture}